\documentclass[a4paper,10pt]{report}\input{../0.Préambule/Préambule_cours_prof}\begin{document}

\newpage
\section*{Périmètre et aires}
\addcontentsline{toc}{section}{Périmètre et aires}

\subsection*{Périmètre et unité}
\addcontentsline{toc}{section}{Périmètre et unité}

\begin{exer1}
Complète.

\begin{enumerate}[font=\color{exer1}]
\renewcommand{\arraystretch}{1.2}
\begin{tabularx}{\linewidth}{m{1.7cm}*{1}{>{\arraybackslash}X}m{0.5cm}m{2cm}*{1}{>{\arraybackslash}X}}
\item $1$km $= $ & \dotfill m&&
\item $1$ cm $= $ & \dotfill mm\\
\item $1$ m $= $ & \dotfill cm&&
\item $1$ dam $= $ & \dotfill m\\
\end{tabularx}
\end{enumerate}
\end{exer1}

\begin{exer1}
Complète.

\begin{enumerate}[font=\color{exer1}]
\renewcommand{\arraystretch}{1.2}
\begin{tabularx}{\linewidth}{m{1.7cm}*{1}{>{\arraybackslash}X}m{0.5cm}m{2cm}*{1}{>{\arraybackslash}X}}
\item $3$km $= $ & \dotfill m&&
\item $5$ mm $= $ & \dotfill m\\
\item $2$ cm $= $ & \dotfill m&&
\item $3$ m $= $ & \dotfill km\\
\item $7$ m $= $ & \dotfill cm&&
\item $5$ cm $= $ & \dotfill dm\\
\end{tabularx}
\end{enumerate}
\end{exer1}

\begin{exer1}
Je suis un rectangle $ABCD$. $[AB]$ mesure $36$mm et $[BC]$ mesure $4$cm.
\begin{enumerate}[font=\color{exer1}]
\item Fais un dessin a main levée
\item Quelle est mon périmètre.
\end{enumerate}
\end{exer1}

\noindent \begin{minipage}[t]{8.3cm}
\begin{exer1}
Donne la nature du quadrilatère en justifiant, puis calcule son périmètre.

\includegraphics[scale=0.5]{../40.Image/17. Exo-Longueur, périmètre et aire 5.png}
\end{exer1}
\end{minipage}
\hspace{5mm}
\begin{minipage}[t]{8.3cm}
\begin{exer1}
Calcule le périmètre (donne à chaque fois la valeur exacte, puis arrondie au dixième près)
\begin{enumerate}[font=\color{exer1}]
\item d'un disque de rayon $4$cm
\item d'un disque de diamètre $12$cm
\end{enumerate}
\end{exer1}
\end{minipage}

\subsection*{Problèmes}
\addcontentsline{toc}{section}{Problèmes}

\vspace{-5mm}
\noindent \begin{minipage}{7.5cm}
\begin{exer3}
Sachant que $AB=9$cm ; $BC=21$cm ; $CD = 11$cm ; $DE = 9$cm ; $EF = 11$cm ; $GH = 7$cm; calcule le périmètre de rectangle $ACEG$

\begin{center}
\includegraphics[scale=0.5]{../40.Image/17. Exo-Longueur, périmètre et aire 10.png}
\end{center}
\end{exer3}

\begin{exer3}
On considère un rectangle de largeur $l$ et de longueur $L$. On triple sa longueur et on triple sa largeur. Par combien est multiplié le périmètre de ce rectangle ?
\end{exer3}
\end{minipage}
\hspace{5mm}
\begin{minipage}{9cm}
\begin{exer3}[Parcours de santé]

\begin{enumerate}[font=\color{exer3}]
\item Sur le parcours de santé ci-dessous, repasse en vert les petits arc de cercle de même rayon et repasse en rouge les grands arcs de cercle de même rayon.

\begin{center}
\includegraphics[scale=0.5]{../40.Image/17. Exo-Longueur, périmètre et aire 11.png}
\end{center}
\item Calcule la longueur réelle du parcours de santé, au mètre près.
\end{enumerate}
\end{exer3}
\end{minipage}

\subsection*{Conversion d'unités d'aires}
\addcontentsline{toc}{section}{Conversion d'unités d'aires}

\begin{exer1}
Complète.

\begin{enumerate}[font=\color{exer1}]
\renewcommand{\arraystretch}{1.2}
\begin{tabularx}{\linewidth}{m{1.7cm}*{1}{>{\arraybackslash}X}m{0.5cm}m{2cm}*{1}{>{\arraybackslash}X}}
\item $3$m$^2$ $= $ & \dotfill cm$^2$&&
\item $700$cm$^2$ $= $ & \dotfill m$^2$\\
\item $5$m$^2$ $= $ & \dotfill dm$^2$&&
\item $6000$dm$^2$ $= $ & \dotfill m$^2$\\
\item $5$cm$^2$ $= $ & \dotfill mm$^2$&&
\item $400$mm$^2$ $= $ & \dotfill cm$^2$\\
\end{tabularx}
\end{enumerate}
\end{exer1}

\begin{exer2}
On donne les superficies suivantes :
\begin{enumerate}[font=\color{exer2}]
\renewcommand{\arraystretch}{0.7}
\begin{tabularx}{\linewidth}{*{2}{>{\arraybackslash}X}}
\item[•] Les Bermudes : $53$km$^2$&
\item[•] Iles Tuvalu : $\np{2600}$ ha \\
\item[•] Ile Christmas : $\np{135000000}$m$^2$&
\item[•] Saint-Pierre-et-Miquelon	
 : $\np{2430000}$dam$^2$\\
\end{tabularx}
\end{enumerate}
Range ces iles dans l'ordre décroissant de leur superficie.
\end{exer2}

\subsection*{Calculer une aire}
\addcontentsline{toc}{section}{Calculer une aire}

\vspace{-1cm}
\noindent \begin{minipage}{10cm}
\begin{exer1}
Donne la nature de chaque quadrilatère en justifiant, puis calcule (sans calculatrice) son aire (pose si nécessaire les opérations)
\end{exer1}
\end{minipage}
\hspace{3mm}
\begin{minipage}{6cm}
\begin{center}
\includegraphics[scale=0.6]{../40.Image/14. Exo-Les aires 7.png}
\end{center}
\end{minipage}

\begin{exer1}
Calcule les aires de chacun des triangles. (Attention, les triangles ne sont pas dessinés en vraie grandeur.)

\begin{center}
\includegraphics[scale=0.65]{../40.Image/14. Exo-Les aires 10.png}
\end{center}
\end{exer1}

\begin{exer2}
$HSLE$ est un rectangle d'aire $18$cm$^2$ et $L=6cm$. Détermine la longueur du segment $[LS]$.
\end{exer2}

\begin{exer2}
$RSV$ est un triangle rectangle en $S$ d'aire $7$cm$^2$ et $RS=2cm$. Détermine la longueur de $[VS]$.
\end{exer2}

\begin{exer2}
$GSLM$ est un carrée. $[GS]$ mesure $3$cm. $GSB$ est un triangle rectangle en $S$, $B$ étant à l'extérieur du carré. $[BS]$ mesure $5$cm. Trace une figure a main levée, codée, puis calcule sont aire totale.
\end{exer2}

\begin{exer1}
Calcule l'aire d'un disque de rayon $4$cm et l'aire d'un disque de diamètre $12$cm. On donnera le calcul (avec $\pi$) et le résultat arrondie au dixième près.
\end{exer1}

\begin{exer1}
Calcule l'aire d'un disque de rayon $4$m et l'aire d'un disque de diamètre $4,3$hm. On donnera le calcul (avec $\pi$) et le résultat arrondie au centième près.
\end{exer1}

\noindent \begin{minipage}{9cm}
\begin{exer1}
Calcule l'aire des figures colorées suivantes. Donne le calcul (avec le nombre $\pi$), puis la valeur arrondie au centième.
\end{exer1}
\end{minipage}
\begin{minipage}{8cm}
\begin{center}
\includegraphics[scale=0.65]{../40.Image/14. Exo-Les aires 12.png}
\end{center}
\end{minipage}
\end{document}